\documentclass[12pt,a4paper,oneside]{article}
\usepackage[brazil]{babel}
\usepackage[utf8]{inputenc}
\usepackage{olymp}

\usepackage{wrapfig}

\setcounter{problem}{1}

\begin{document}

\begin{problem}{Bola de boliche}{bola}{2}
 
\begin{wrapfigure}{r}{0.3\textwidth}
    \includegraphics[width=0.3\textwidth]{images/exercise_5_0.jpg}
\end{wrapfigure}

A Copa do Qatar 2022 foi realizada no final do ano passado, tendo a Al Rihla, da foto ao lado, como a bola oficial. Bolas de futebol são muito fáceis de transportar, já que elas saem das fábricas vazias e são enchidas somente pelas lojas ou pelos consumidores finais. Infelizmente o mesmo não pode ser dito das bolas de boliche. Como são completamente sólidas, elas só podem ser transportadas embaladas uma a uma, em caixas separadas.

Certa fábrica de bolas de boliche trabalha somente através de encomendas e envia todas as bolas por SEDEX. Como as bolas têm tamanhos diferentes, a fábrica tem vários tamanhos diferentes de caixas para transportá-las.

Escreva um programa que, dado o diâmetro de uma bola e as 3 dimensões de uma caixa (altura, largura e profundidade), diz se a bola de boliche cabe dentro da caixa ou não.

%\Notes
%
%Observações, se tiver.

\InputFile

A primeira linha da entrada contém um inteiro $D$ ($1 \leq D \leq 10000$) que indica o diâmetro da bola de boliche. A segunda linha da entrada contém 3 números inteiros separados por um espaço cada: a altura $A$ ($1 \leq A \leq 10000$), seguida da largura $L$ ($1 \leq L \leq 10000$) e da profundidade $P$ ($1 \leq P \leq 10000$).

%\Notes



\OutputFile

Seu programa deve imprimir uma única linha, contendo a letra ‘\texttt{S}’ caso a bola de boliche caiba dentro da caixa ou ‘\texttt{N}’ caso contrário.

\Example

\begin{example}
\exmp{
3
2 3 5
}{
N
}%
\end{example}

\begin{example}
\exmp{
5
5 5 5
}{
S
}%
\end{example}

\begin{example}
\exmp{
9
15 9 10
}{
S
}%
\end{example}


\footnotesize Problema da OBI 2010 - Olimpíada Brasileira de Informática (era Copa de 2010 na África do Sul...)

\end{problem}

\end{document}
